\section{Project Types and Phases}


\begin{frame}{Projects}
    \mytext{
    
    \begin{itemize}%[<+->]
        \item What is \emph{not }a project? $\rightarrow$ `business as usual'
        \item Projects have a beginning, a middle and an end...
        \item Someone \emph{wants }a project
        \item Project \emph{makes a difference} $\rightarrow$ change
        \item Projects have specifics
        \item Projects involve people and other resources
        \item Projects have sponsors, owners, managers...
        \item Projects can be small, large...
        \item Projects can be services, actions, products...
        \item Products can be tangible, intangible...
        \item Projects are planned...
    \end{itemize}
    }

\end{frame}

\begin{frame}{Example projects}
    \mytext{
    Think of one project example!
    }
\end{frame}
\begin{frame}{Example projects}
    \mytext{
    \begin{itemize}%[<+->]
        \item Office move
        \item Construction project
        \item Business expanse
        \item Business acquisition
        \item Conference, event
        \item IT project
        \item Moon landing
        \item Military projects
        \item Privatisation
        \item Founding of an institution
        \item Disbanding, dismanteling
    \end{itemize}
    }
\end{frame}

\begin{frame}{Projects are \textbf{not}...}
    \mytext{
    
        \begin{itemize}
            \item wishes, ideas or ideals 
            \item vague goals or directions
            \item positions or functions
        \end{itemize}
    
        \medskip
    
        $\rightarrow$ projects need specifics 
    }
\end{frame}    

\begin{frame}{Projects are \textbf{not}...}
    \mytext{
     
        {\LARGE Products $\neq$ PROJECTS!}
    
        \medskip
    
        \begin{itemize}
            \item The anti-malaria product
            \item Mr. Smith and his qualities
            \item The presidency of Mr.\ Smith
            \item ...
        \end{itemize}
    
        A project is \textbf{a planned process,} and \textbf{its result} is a product!
    
        \medskip
    
        So: how to get from the \emph{projected} result to the \emph{project?} 
    }
\end{frame}





\begin{frame}{Definition PMI}
    \mytext{
    
    \begin{quote}
        % \singlespacing
        \emph{All projects are a temporary effort to create value through a unique product, service or result. All projects have a beginning and an end. They have a team, a budget, a schedule and a set of expectations the team needs to meet. Each project is unique and differs from routine operations—the ongoing activities of an organization—because projects reach a conclusion once the goal is achieved.}
        
        \emph{The changing nature of work due to technological advances, globalization and other factors means that, increasingly, work is organized around projects with teams being brought together based on the skills needed for specific tasks.}
        \par
        \hfill ---Project Management Institute \cite{PMIWeb}
    \end{quote}
    }
\end{frame}

\begin{frame}{Definition PMI}

    \mytext{
    \begin{itemize}
        \item effort to create value
        \item through a unique project, service or result
        \item beginning and end
        \item team
        \item budget
        \item schedule
        \item set of expectations
        \item closure; $\neq$ routine operations
    \end{itemize}
    }
\end{frame}
    
\begin{frame}{Definition PMI}
    \mytext{
    I like to add:
    
    \begin{itemize}
        \item \emph{work breakdown / sub-projects}
        \item \emph{team (management)}
        \item \emph{stakeholders and contractors}
        \item \emph{governance and legal structure}
        \item \emph{technique and culture}
        \item \emph{risks / opportunities / alternatives}
        \item \emph{external effects in space and time}
        \item \emph{standards and regulations}
    \end{itemize}
    }
\end{frame}

\begin{frame}{Qualities of Projects (non-exahustive list)}
\mytext{

    \begin{itemize}
        \item Large or small
        \item Complex or simple
        \item Environment: dynamic, static
        \item Clearly defined, loosely defined
        \item New or established
        \item Initialised, projected
        \item Global, local
        \item Urgent, not urgent
        \item Risky, safe
        \item High priority, lower priority
    \end{itemize}
}
\end{frame}

\begin{frame}{Types of Projects - by outcome}
    \showmytikz{../shared/tikz/forest-projecttypes}
\end{frame}


\begin{frame}{Types of Projects - by owner}
    \showmytikz{../shared/tikz/forest-projects-by-owner}
\end{frame}

\begin{frame}{Ownership Hierarchy}
    \showmytikz{../shared/tikz/tikz-ownership-hierarchy}
    % \begin{tikzpicture}
    % \tikzset{every node/.style={minimum height=3ex,draw,rounded corners}}
    %     \node  (sponsor) at (0,0) {Project Sponsor};
    %     \node[right=of sponsor] (public) {Public Bodies};
    %     \node[below of=sponsor] (holder)   {Project Holder (\emph{Projektträger})};
    %     \node[right=of holder] (consultants) {Consultants};
    %     \node[below of=holder] (pmo)  {Project Management Office(s)};
    %     \node[below of=pmo] (contractor)  {Contractors};
    %     \node[right=of pmo] (others) {...};
    %     \node[right=of contractor] (users) {(Future) Users};
    %     \draw[->] (sponsor) -- (holder);
    %     \draw[->]  (holder) -- (pmo);
    %     \draw[->] (pmo) -- (contractor);
    % \end{tikzpicture}
\end{frame}

\begin{frame}{The Three Constraints}
    \showmytikz[caption={The Three Constraints (Introduction)},label=fig:threeconstraints-intro]{../shared/tikz/tikz-threeconstraints}

\end{frame}



\begin{frame}{Types of Projects - by size}
    \showmytikz{../shared/tikz/tikz-projects-by-size}
\end{frame}



\begin{frame}{Sub-Projects: Work Breakdown Structure}
    \showmytikz{../shared/tikz/forest-projects-sizes-breakdown}
\end{frame}

\begin{frame}{Sub-Projects: Work Breakdown Structure}
    \showmytikz{../shared/tikz/forest-work-packages}
\end{frame}

\begin{frame}{Figures / Numbers: KPIs}
    \mytext{
    
    \begin{itemize}
        \item KPI Key Performance Indicator
        \item Money, time, quality, quantities...
        \item "What you can measure, you can manage"
        \item Data-driven decisions
    \end{itemize}
    
    \medskip

    When talking about anything in business,\\ provide figures / numbers!
    }
\end{frame}

\begin{frame}{Orders of Magnitude}
\mytable{

    \begin{tabular}{lrrll}
        \toprule
        \multicolumn{2}{l}{science} & \multicolumn{2}{l}{business english} & \textit{avoid}\\
        \midrule
         $10^3$ & K & T & thousand & \\
        $10^4$ & 10K & T & & \\
        $10^5$ & 100K & T & & \emph{lakh}\\
        \midrule
        $10^6$ & 1M & M & million & \emph{Millionen} \\
        $10^7$ & 10M & M & &  \emph{crore}\\
        $10^8$ & 100M & M & & \\
        \midrule
        $10^9$ & 1G & B & billion & \emph{Milliarden} \\
        $10^{10}$ & 10G & B& & \\
        $10^{11}$ & 100G & B & & \\
        \midrule
        $10^{12}$ & 1T &  & trillion & \emph{Billionen}\\
        \bottomrule
    \end{tabular}
}

\end{frame}

\begin{frame}{Types of Projects - by budget}
\mytext{

    \mytable{
    
        \begin{tabular}{rl}
        \toprule
        bracket & \textit{example} \\
        \midrule
            10--50 TEUR & \\
            50--100 TEUR & \\
            100--250 TEUR & \\
            250--500 TEUR & \\
            500--1000 TEUR & \\
            1--2 MEUR & \\
            2--5 MEUR & \\
            5--10 MEUR & \\
            ... & \\
        \bottomrule
        \end{tabular}
    }

    \bigskip

    Loft bed? Drilling platform? Airport?

    Bridge? Ferry boat? Office move?

    SAP introduction? Marketing campaign?
    
}

\end{frame}

\begin{frame}{(Potential) Project Key Data}
    \mytext{
    
    \begin{itemize}
        \item Budget: pre/post
        \begin{itemize}
            \item typically underestimated
            \item often by a factor of 2-5
            \item often by an order of magnitude
        \end{itemize}
        \item Budget variance?
        \item Schedule: duration etc.
        \item Team, Stakeholders
        \item Complexity measures
        \item Sizes, weights, measures
        \item All sorts of business key figure
    \end{itemize}
    }
\end{frame}


\begin{frame}{Project Phases, Ideally}
    \showmytikz{../shared/tikz/tikz-chain-projectphases-ideal}
\end{frame}

\begin{frame}{Project Phases, Realistically}
    \showmytikz{../shared/tikz/tikz-chain-projectphases-morereal}
\end{frame}

\begin{frame}{Project Phases per Schliersmann}
    \showmytikz{../shared/tikz/tikz-phase-oriented-project-management}
    \mycomment{
    
    
    \cite{schiersmannthiel}
    }
\end{frame}

\begin{frame}{PDCA-Cylce}
    \showmytikz[caption={PDCA Cycle (Introduction)},label=fig:pdca-cycle-intro]{../shared/tikz/tikz-PM-smartd-pdca-cycle}
\end{frame}



\begin{frame}{Project Phases (Roberts)}
    \showmytikz[1]{../shared/tikz/tikz-chain-projectphases-details}
\end{frame}

\begin{frame}{Project Phases (Roberts)}
    \showmytikz[1,2]{../shared/tikz/tikz-chain-projectphases-details}
\end{frame}

\begin{frame}{Project Phases (Roberts)}
    \showmytikz[1,2,3]{../shared/tikz/tikz-chain-projectphases-details}
\end{frame}

\begin{frame}{Project Phases (Roberts)}
    \showmytikz[1,2,3,4]{../shared/tikz/tikz-chain-projectphases-details}
\end{frame}

\begin{frame}{Project Phases (Roberts)}
    \showmytikz[1,2,3,4,5]{../shared/tikz/tikz-chain-projectphases-details}
\end{frame}

\begin{frame}{Project Phases (Roberts)}
    \showmytikz[4,5,6]{../shared/tikz/tikz-chain-projectphases-details}
\end{frame}

\begin{frame}{Project Phases (Roberts)}
    \showmytikz[4,5,6,7]{../shared/tikz/tikz-chain-projectphases-details}
\end{frame}

\begin{frame}{Project Phases (Roberts)}
    \showmytikz[1,2,3,4,5,6,7]{../shared/tikz/tikz-chain-projectphases-details}
\end{frame}



\begin{frame}{The Main Phases / Preliminary Gantt}
    \resizebox{\textwidth}{!}{
    \begin{meinganttq}{5}
        \inserttable{gantts/mains.csv}
        \showmilestones{gantts/mains.csv}
    \end{meinganttq}
    }
\end{frame}

\begin{frame}{Cost Components over time}
        \showmytikz{../shared/tikz/tikz-costsovertime.tex}
\end{frame}




\begin{frame}{Project Governance Report}
    \mytext{
    
    \begin{itemize}
        \item The `constitution' of a project
        \item A contract between stakeholders
        \item The actual plan, typically comprising
        \begin{itemize}
            \item user requirements ('need' or reason)
            \item detailed solution design
            \item stakeholder analysis
            \item work breakdown structure 
            \item definition of work packages and responsibilities
            \item project plan: network / BPMN plan
            \item project plan: temporal plan / Gantt diagram
            \item governance plan: corporation, controls
            \item identity and roles
            \item communication strategy 
            \item documentation, controls and reporting
            \item risk register
        \end{itemize}
    \end{itemize}
    }
\end{frame}



